\documentclass{article}
\usepackage{ctex} % 引入中文支持
\usepackage{multirow} % 用于合并行
\usepackage{makecell} % 用于单元格内强制换行实现竖排文字
\usepackage{geometry}
\geometry{a4paper, landscape, margin=1in} % 页面横向以适应宽表格

\begin{document}

\begin{table}[htbp]
    \centering
    \renewcommand{\arraystretch}{1.5} % 增加行高，让表格显得不那么拥挤
    \begin{tabular}{|c|c|c|c|c|c|c|c|c|c|c|c|c|}
        \hline
        % 第一行
        & \multicolumn{6}{c|}{整数部分} 
        & 小数点 
        & \multicolumn{5}{c|}{小数部分} \\
        \hline
        
        % 第二行
        数位 
        & $\cdots$ 
        & \makecell{万\\位} 
        & \makecell{千\\位} 
        & \makecell{百\\位} 
        & \makecell{十\\位} 
        & \makecell{个\\位} 
        & \multirow{2}{*}{$\cdot$} % 合并两行作为小数点
        & \makecell{十\\分\\位} 
        & % 故意留空的百分位
        & \makecell{千\\分\\位} 
        & \makecell{万\\分\\位} 
        & $\cdots$ \\
        \cline{1-7} \cline{9-13} % 画横线，跳过第8列（小数点列）
        
        % 第三行
        \makecell{计\\数\\单\\位} 
        & $\cdots$ 
        & 万 & 千 & 百 & 十 
        & \makecell{一\\(个)} 
        & % 此处已被 multirow 占用
        & \makecell{十\\分\\之\\一} 
        & \makecell{百\\分\\之\\一} 
        & \makecell{千\\分\\之\\一} 
        & % 故意留空的万分之一
        & $\cdots$ \\
        \hline
    \end{tabular}
\end{table}

\end{document}